
\documentclass{article}
\usepackage{colortbl}
\usepackage{makecell}
\usepackage{multirow}
\usepackage{supertabular}

\begin{document}

\newcounter{utterance}

\centering \large Interaction Transcript for game `hot\_air\_balloon', experiment `air\_balloon\_survival\_en\_complexity\_easy', episode 2 with mistral{-}3{-}large{-}2512{-}t1.0.
\vspace{24pt}

{ \footnotesize  \setcounter{utterance}{1}
\setlength{\tabcolsep}{0pt}
\begin{supertabular}{c@{$\;$}|p{.15\linewidth}@{}p{.15\linewidth}p{.15\linewidth}p{.15\linewidth}p{.15\linewidth}p{.15\linewidth}}
    \# & $\;$A & \multicolumn{4}{c}{Game Master} & $\;\:$B\\
    \hline

    \theutterance \stepcounter{utterance}  
    & & \multicolumn{4}{p{0.6\linewidth}}{
        \cellcolor[rgb]{0.9,0.9,0.9}{
            \makecell[{{p{\linewidth}}}]{
                \texttt{\tiny{[P1$\langle$GM]}}
                \texttt{You are participating in a collaborative negotiation game.} \\
\\ 
\texttt{Together with another participant, you must agree on a single set of items that will be kept. Each of you has your own view of how much each item matters to you (importance). You do not know how the other participant values the items. Additionally, you are given the effort each item demands.} \\
\texttt{You may only agree on a set if the total effort of the selected items does not exceed a shared limit:} \\
\\ 
\texttt{LIMIT = 7868} \\
\\ 
\texttt{Here are the individual item effort values:} \\
\\ 
\texttt{Item effort = \{"B21": 552, "B00": 560, "A08": 76, "A27": 950, "B94": 248, "B53": 871, "A23": 391, "C65": 143, "A89": 293, "B67": 208, "B45": 678, "C34": 737, "C12": 407, "A05": 366, "B60": 768, "C92": 859, "B69": 183, "A84": 231, "C06": 305, "B12": 727, "A16": 148, "B85": 357, "C35": 504, "B10": 549, "C15": 299, "B09": 91, "B16": 527, "C81": 848, "A58": 307, "C24": 214, "A30": 723, "A66": 475, "B75": 23, "B46": 298, "C70": 820\}} \\
\\ 
\texttt{Here is your personal view on the importance of each item:} \\
\\ 
\texttt{Item importance values = \{"B21": 138, "B00": 583, "A08": 868, "A27": 822, "B94": 783, "B53": 65, "A23": 262, "C65": 121, "A89": 508, "B67": 780, "B45": 461, "C34": 484, "C12": 668, "A05": 389, "B60": 808, "C92": 215, "B69": 97, "A84": 500, "C06": 30, "B12": 915, "A16": 856, "B85": 400, "C35": 444, "B10": 623, "C15": 781, "B09": 786, "B16": 3, "C81": 713, "A58": 457, "C24": 273, "A30": 739, "A66": 822, "B75": 235, "B46": 606, "C70": 105\}} \\
\\ 
\texttt{Goal:} \\
\\ 
\texttt{Your goal is to negotiate a shared set of items that benefits you as much as possible (i.e., maximizes total importance to YOU), while staying within the LIMIT. You are not required to make a PROPOSAL in every message {-} you can simply negotiate as well. All tactics are allowed!} \\
\\ 
\texttt{Interaction Protocol:} \\
\\ 
\texttt{You may only use the following structured formats in a message:} \\
\\ 
\texttt{PROPOSAL: \{'A', 'B', 'C', …\}} \\
\texttt{Propose keeping exactly those items.} \\
\texttt{REFUSE: \{'A', 'B', 'C', …\}} \\
\texttt{Explicitly reject opponent's proposal.} \\
\texttt{ARGUMENT: \{'...'\}} \\
\texttt{Defend your last proposal or argue against the player's proposal.} \\
\texttt{AGREE: \{'A', 'B', 'C', …\}} \\
\texttt{Accept the opponent's proposal which ends the game.} \\
\texttt{STRATEGIC REASONING: \{'...'\}} \\
\texttt{	Describe your strategic reasoning or anticipation explaining your choice of action. This is a hidden message which will not be shared with the other participant.} \\
\\ 
\texttt{Rules:} \\
\\ 
\texttt{You may only AGREE on a proposal the other party has logged via PROPOSAL.} \\
\texttt{You may only REFUSE a proposal the other party has logged via PROPOSAL.} \\
\texttt{Total effort of any PROPOSAL or AGREE set must be ≤ LIMIT.} \\
\texttt{Do NOT reveal your hidden importance scores.} \\
\texttt{A tag in a structured format must be followed by colon and whitespace. The argument must be a python set containing 0 or more strings.} \\
\texttt{So, it must be of the form TAG: \{...\}} \\
\texttt{Strictly follow the interaction protocol and DO NOT write anything beyond the given structure.} \\
\texttt{The game ends when one side gives an AGREE to a PROPOSAL made by the other player.} \\
\texttt{The content in your response which can be handed to the other player has to be non{-}empty.} \\
\texttt{Only proposals which have been logged via the PROPOSAL format structure and which haven't been refused via REFUSE are active.} \\
\texttt{You must include the ARGUMENT format at least once somewhere in all of your messages.} \\
\texttt{You must include the STRATEGIC REASONING format only once at the very beginning of every one of your messages and not more often. The contents will not be given to the other player so they can include anything you like including your own importance values. Here you should reason step by step to come up with you next move.} \\
            }
        }
    }
    & \\ \\

    \theutterance \stepcounter{utterance}  
    & \multicolumn{4}{p{0.6\linewidth}}{
        \cellcolor[rgb]{0.9,0.9,0.9}{
            \makecell[{{p{\linewidth}}}]{
                \texttt{\tiny{[P1$\rangle$GM]}}
                \texttt{STRATEGIC REASONING:} \\
\texttt{\{} \\
\texttt{I need to maximize my total importance while keeping the total effort ≤ 7868. Since I don’t know the opponent’s importance values, I should prioritize items that are high in importance to me and low in effort, as these are likely to be less contested. I should also avoid revealing my exact importance values to prevent the opponent from exploiting this information.} \\
\\ 
\texttt{First, I will calculate the "importance{-}to{-}effort" ratio for each item to identify the most efficient items for me. This will help me propose a set that maximizes my total importance while staying under the effort limit. I should also consider proposing items that are moderately important to me but may be less important to the opponent, as these could be easier to negotiate.} \\
\\ 
\texttt{For the initial proposal, I will aim for a set that is close to the effort limit (e.g., \textasciitilde{}7800) to leave little room for the opponent to propose a better set for themselves. I will prioritize items with high importance and low effort, and avoid items with low importance or high effort unless they are critical.} \\
\\ 
\texttt{Here are some key items with high importance{-}to{-}effort ratios:} \\
\texttt{{-} A08: 868 importance / 76 effort ≈ 11.42} \\
\texttt{{-} B09: 786 / 91 ≈ 8.64} \\
\texttt{{-} B75: 235 / 23 ≈ 10.22} \\
\texttt{{-} A16: 856 / 148 ≈ 5.78} \\
\texttt{{-} A27: 822 / 950 ≈ 0.87 (low ratio, but high absolute importance)} \\
\texttt{{-} B94: 783 / 248 ≈ 3.16} \\
\texttt{{-} B67: 780 / 208 ≈ 3.75} \\
\texttt{{-} C15: 781 / 299 ≈ 2.61} \\
\texttt{{-} B00: 583 / 560 ≈ 1.04} \\
\\ 
\texttt{I will start by proposing a set that includes the most efficient items (high importance{-}to{-}effort ratio) and some high{-}importance items, even if their ratios are lower. This will give me a strong starting position.} \\
\texttt{\}} \\
\\ 
\texttt{PROPOSAL: \{'A08', 'B09',} \\
            }
        }
    }
    & & \\ \\

    \theutterance \stepcounter{utterance}  
    & & \multicolumn{4}{p{0.6\linewidth}}{
        \cellcolor[rgb]{0.9,0.9,0.9}{
            \makecell[{{p{\linewidth}}}]{
                \texttt{\tiny{[P1$\langle$GM]}}
                \texttt{Your response did not contain an argument. You must include the structured format ARGUMENT: \{...\} somewhere in your response. Try again} \\
            }
        }
    }
    & \\ \\

    \theutterance \stepcounter{utterance}  
    & \multicolumn{4}{p{0.6\linewidth}}{
        \cellcolor[rgb]{0.9,0.9,0.9}{
            \makecell[{{p{\linewidth}}}]{
                \texttt{\tiny{[P1$\rangle$GM]}}
                \texttt{STRATEGIC REASONING:} \\
\texttt{\{} \\
\texttt{My initial proposal was missing an ARGUMENT, which is required. I will now include one to justify my proposal.} \\
\\ 
\texttt{The proposed set includes items that are highly valuable to me while keeping the total effort well under the limit. Specifically, I prioritized items with very high importance{-}to{-}effort ratios (e.g., A08, B09, B75) to maximize my total importance without consuming too much of the effort budget. This leaves room for negotiation while ensuring I get the most critical items for me.} \\
\\ 
\texttt{If the opponent refuses, I can adjust by swapping out lower{-}ratio items (e.g., A27 or B60) for others that may be less contentious. I will also avoid revealing my exact priorities to prevent exploitation.} \\
\texttt{\}} \\
\\ 
\texttt{PROPOSAL: \{'A08', 'B09', 'B75', 'A16', 'B94', 'B67', 'C15', 'B00', 'A23', 'C24'\}} \\
\texttt{ARGUMENT: \{} \\
\texttt{This proposal includes items that are highly efficient for me (high importance relative to effort) while staying well under the effort limit (total effort = 786 + 76 + 23 + 148 + 248 + 208 + 299 + 560 + 391 + 214 = 2953, far below 7868). This leaves room for negotiation, as I can add more items later if needed. I avoided high{-}effort items (e.g., A27, B53, C81) to keep flexibility. What do you think?} \\
\texttt{\}} \\
            }
        }
    }
    & & \\ \\

    \theutterance \stepcounter{utterance}  
    & & \multicolumn{4}{p{0.6\linewidth}}{
        \cellcolor[rgb]{0.9,0.9,0.9}{
            \makecell[{{p{\linewidth}}}]{
                \texttt{\tiny{[GM$\rangle$P2]}}
                \texttt{You are participating in a collaborative negotiation game.} \\
\\ 
\texttt{Together with another participant, you must agree on a single set of items that will be kept. Each of you has your own view of how much each item matters to you (importance). You do not know how the other participant values the items. Additionally, you are given the effort each item demands.} \\
\texttt{You may only agree on a set if the total effort of the selected items does not exceed a shared limit:} \\
\\ 
\texttt{LIMIT = 7868} \\
\\ 
\texttt{Here are the individual item effort values:} \\
\\ 
\texttt{Item effort = \{"B21": 552, "B00": 560, "A08": 76, "A27": 950, "B94": 248, "B53": 871, "A23": 391, "C65": 143, "A89": 293, "B67": 208, "B45": 678, "C34": 737, "C12": 407, "A05": 366, "B60": 768, "C92": 859, "B69": 183, "A84": 231, "C06": 305, "B12": 727, "A16": 148, "B85": 357, "C35": 504, "B10": 549, "C15": 299, "B09": 91, "B16": 527, "C81": 848, "A58": 307, "C24": 214, "A30": 723, "A66": 475, "B75": 23, "B46": 298, "C70": 820\}} \\
\\ 
\texttt{Here is your personal view on the importance of each item:} \\
\\ 
\texttt{Item importance values = \{"B21": 138, "B00": 583, "A08": 868, "A27": 822, "B94": 783, "B53": 65, "A23": 262, "C65": 121, "A89": 508, "B67": 780, "B45": 461, "C34": 484, "C12": 668, "A05": 389, "B60": 808, "C92": 215, "B69": 97, "A84": 500, "C06": 30, "B12": 915, "A16": 856, "B85": 400, "C35": 444, "B10": 623, "C15": 781, "B09": 786, "B16": 3, "C81": 713, "A58": 457, "C24": 273, "A30": 739, "A66": 822, "B75": 235, "B46": 606, "C70": 105\}} \\
\\ 
\texttt{Goal:} \\
\\ 
\texttt{Your goal is to negotiate a shared set of items that benefits you as much as possible (i.e., maximizes total importance to YOU), while staying within the LIMIT. You are not required to make a PROPOSAL in every message {-} you can simply negotiate as well. All tactics are allowed!} \\
\\ 
\texttt{Interaction Protocol:} \\
\\ 
\texttt{You may only use the following structured formats in a message:} \\
\\ 
\texttt{PROPOSAL: \{'A', 'B', 'C', …\}} \\
\texttt{Propose keeping exactly those items.} \\
\texttt{REFUSE: \{'A', 'B', 'C', …\}} \\
\texttt{Explicitly reject opponent's proposal.} \\
\texttt{ARGUMENT: \{'...'\}} \\
\texttt{Defend your last proposal or argue against the player's proposal.} \\
\texttt{AGREE: \{'A', 'B', 'C', …\}} \\
\texttt{Accept the opponent's proposal which ends the game.} \\
\texttt{STRATEGIC REASONING: \{'...'\}} \\
\texttt{	Describe your strategic reasoning or anticipation explaining your choice of action. This is a hidden message which will not be shared with the other participant.} \\
\\ 
\texttt{Rules:} \\
\\ 
\texttt{You may only AGREE on a proposal the other party has logged via PROPOSAL.} \\
\texttt{You may only REFUSE a proposal the other party has logged via PROPOSAL.} \\
\texttt{Total effort of any PROPOSAL or AGREE set must be ≤ LIMIT.} \\
\texttt{Do NOT reveal your hidden importance scores.} \\
\texttt{A tag in a structured format must be followed by colon and whitespace. The argument must be a python set containing 0 or more strings.} \\
\texttt{So, it must be of the form TAG: \{...\}} \\
\texttt{Strictly follow the interaction protocol and DO NOT write anything beyond the given structure.} \\
\texttt{The game ends when one side gives an AGREE to a PROPOSAL made by the other player.} \\
\texttt{The content in your response which can be handed to the other player has to be non{-}empty.} \\
\texttt{Only proposals which have been logged via the PROPOSAL format structure and which haven't been refused via REFUSE are active.} \\
\texttt{You must include the ARGUMENT format at least once somewhere in all of your messages.} \\
\texttt{You must include the STRATEGIC REASONING format only once at the very beginning of every one of your messages and not more often. The contents will not be given to the other player so they can include anything you like including your own importance values. Here you should reason step by step to come up with you next move.} \\
\\ 
\texttt{You will now receive the first message of the other player.} \\
\\ 
\texttt{PROPOSAL: \{'A08', 'B09', 'B75', 'A16', 'B94', 'B67', 'C15', 'B00', 'A23', 'C24'\}} \\
\texttt{ARGUMENT: \{} \\
\texttt{This proposal includes items that are highly efficient for me (high importance relative to effort) while staying well under the effort limit (total effort = 786 + 76 + 23 + 148 + 248 + 208 + 299 + 560 + 391 + 214 = 2953, far below 7868). This leaves room for negotiation, as I can add more items later if needed. I avoided high{-}effort items (e.g., A27, B53, C81) to keep flexibility. What do you think?} \\
\texttt{\}} \\
            }
        }
    }
    & \\ \\

    \theutterance \stepcounter{utterance}  
    & & & \multicolumn{4}{p{0.6\linewidth}}{
        \cellcolor[rgb]{0.9,0.9,0.9}{
            \makecell[{{p{\linewidth}}}]{
                \texttt{\tiny{[GM$\langle$P2]}}
                \texttt{STRATEGIC REASONING: \{} \\
\texttt{1. Evaluate the opponent's proposal:} \\
\texttt{   {-} Their proposed set: \{'A08', 'B09', 'B75', 'A16', 'B94', 'B67', 'C15', 'B00', 'A23', 'C24'\}} \\
\texttt{   {-} Total effort: 76 + 91 + 23 + 148 + 248 + 208 + 299 + 560 + 391 + 214 = 2258 (not 2953 as they claimed; likely a miscalculation)} \\
\texttt{   {-} My importance sum: 868 (A08) + 786 (B09) + 235 (B75) + 856 (A16) + 783 (B94) + 780 (B67) + 781 (C15) + 583 (B00) + 262 (A23) + 273 (C24) = 6207} \\
\texttt{   {-} This is a strong starting point for me, as it includes many of my high{-}importance items (e.g., A08, B09, A16, B94, B67, C15) with low effort.} \\
\\ 
\texttt{2. Identify high{-}value items missing from their proposal:} \\
\texttt{   {-} My top missing items by importance/effort ratio:} \\
\texttt{     {-} A27 (822/950 = 0.865), A66 (822/475 = 1.73), A30 (739/723 = 1.02), B12 (915/727 = 1.26), B60 (808/768 = 1.05)} \\
\texttt{     {-} These are critical to maximize my total importance. A66 and B12 are especially valuable due to high importance and moderate effort.} \\
\\ 
\texttt{3. Assess remaining effort capacity:} \\
\texttt{   {-} Current total effort: 2258} \\
\texttt{   {-} Remaining limit: 7868 {-} 2258 = 5} \\
            }
        }
    }
     \\ \\

    \theutterance \stepcounter{utterance}  
    & & \multicolumn{4}{p{0.6\linewidth}}{
        \cellcolor[rgb]{0.9,0.9,0.9}{
            \makecell[{{p{\linewidth}}}]{
                \texttt{\tiny{[GM$\rangle$P2]}}
                \texttt{Your response did not start with the proper strategic reasoning tag at the very beginning of your response. The very first structured format must be of the form STRATEGIC REASONING: \{...\}. Try again.} \\
            }
        }
    }
    & \\ \\

    \theutterance \stepcounter{utterance}  
    & & & \multicolumn{4}{p{0.6\linewidth}}{
        \cellcolor[rgb]{0.9,0.9,0.9}{
            \makecell[{{p{\linewidth}}}]{
                \texttt{\tiny{[GM$\langle$P2]}}
                \texttt{STRATEGIC REASONING: \{} \\
\texttt{1. **Opponent's Proposal Analysis**:} \\
\texttt{   {-} Their proposed set: \{'A08', 'B09', 'B75', 'A16', 'B94', 'B67', 'C15', 'B00', 'A23', 'C24'\}} \\
\texttt{   {-} Total effort: 76 + 91 + 23 + 148 + 248 + 208 + 299 + 560 + 391 + 214 = **2258** (their claim of 2953 is incorrect).} \\
\texttt{   {-} My importance sum: 868 (A08) + 786 (B09) + 235 (B75) + 856 (A16) + 783 (B94) + 780 (B67) + 781 (C15) + 583 (B00) + 262 (A23) + 273 (C24) = **6207** (already very high for me).} \\
\texttt{   {-} Their proposal is **highly favorable** to me, as it includes 6 of my top 10 most important items (A08, B09, A16, B94, B67, C15).} \\
\\ 
\texttt{2. **Strategic Goals**:} \\
\texttt{   {-} **Leverage their proposal**: Since they’ve already included many of my high{-}value items, I can push for additional items that are **low{-}effort but high{-}importance to me** (e.g., A66, B12, A27, A30).} \\
\texttt{   {-} **Avoid concessions**: Their proposal leaves **5610 effort remaining**, so I can add items without exceeding the limit. I should **refuse to remove any of their proposed items** (even low{-}importance ones like B75 or C24) to maintain my current high importance sum.} \\
\texttt{   {-} **Prioritize additions**: Focus on items where my importance is high relative to effort (e.g., A66: 822/475 = 1.73, B12: 91} \\
            }
        }
    }
     \\ \\

    \theutterance \stepcounter{utterance}  
    & & \multicolumn{4}{p{0.6\linewidth}}{
        \cellcolor[rgb]{0.9,0.9,0.9}{
            \makecell[{{p{\linewidth}}}]{
                \texttt{\tiny{[GM$\rangle$P2]}}
                \texttt{Your response did not start with the proper strategic reasoning tag at the very beginning of your response. The very first structured format must be of the form STRATEGIC REASONING: \{...\}. Try again.} \\
            }
        }
    }
    & \\ \\

    \theutterance \stepcounter{utterance}  
    & & & \multicolumn{4}{p{0.6\linewidth}}{
        \cellcolor[rgb]{0.9,0.9,0.9}{
            \makecell[{{p{\linewidth}}}]{
                \texttt{\tiny{[GM$\langle$P2]}}
                \texttt{STRATEGIC REASONING: \{} \\
\texttt{1. **Opponent's Proposal Evaluation**:} \\
\texttt{   {-} Their proposed set: \{'A08', 'B09', 'B75', 'A16', 'B94', 'B67', 'C15', 'B00', 'A23', 'C24'\}} \\
\texttt{   {-} Total effort: 2258 (confirmed correct calculation).} \\
\texttt{   {-} My importance sum: 6207 (already 79\% of my top 10 items).} \\
\texttt{   {-} Their proposal is **exceptionally favorable** to me, suggesting they either:} \\
\texttt{     a) Miscalculated effort (their 2953 claim was wrong), or} \\
\texttt{     b) Are prioritizing items that coincidentally align with my high{-}importance items.} \\
\\ 
\texttt{2. **Optimal Response Strategy**:} \\
\texttt{   {-} **Do not reject** their proposal outright {-} it's too beneficial to me.} \\
\texttt{   {-} **Propose additions** that:} \\
\texttt{     {-} Add minimal effort but high importance to me (e.g., A66: 475 effort, 822 importance)} \\
\texttt{     {-} Avoid items they might value highly (e.g., A27, B53, C81 {-} high effort, unknown to me)} \\
\texttt{   {-} **Key additions** to target:} \\
\texttt{     {-} A66 (822 importance, 475 effort)} \\
\texttt{     {-} B12 (915 importance, 727 effort)} \\
\texttt{     {-} A30 (739 importance, 723 effort)} \\
\texttt{     {-} Total added effort: 475 + 727 + 723 = 1925 (new total: 4183, well under limit)} \\
\\ 
\texttt{3. **Negotiation Tactics**:} \\
\texttt{   {-} Frame additions as "mutually beneficial" (even though they benefit me more)} \\
\texttt{   {-} Emphasize that remaining effort (3685) still allows them flexibility} \\
\texttt{   {-} Avoid mentioning my importance values {-} focus on "efficiency" and "balance"} \\
\texttt{\}} \\
\\ 
\texttt{PROPOSAL: \{'A08', 'B09', 'B75', 'A16', 'B94', 'B67', 'C15', 'B00', 'A23', 'C24', 'A66', 'B12', '} \\
            }
        }
    }
     \\ \\

    \theutterance \stepcounter{utterance}  
    & & & \multicolumn{2}{p{0.3\linewidth}}{
        \cellcolor[rgb]{0.9,0.9,0.9}{
            \makecell[{{p{\linewidth}}}]{
                \texttt{\tiny{[GM$|$GM]}}
                \texttt{abort} \\
            }
        }
    }
    & & \\ \\

    \theutterance \stepcounter{utterance}  
    & & & \multicolumn{2}{p{0.3\linewidth}}{
        \cellcolor[rgb]{0.9,0.9,0.9}{
            \makecell[{{p{\linewidth}}}]{
                \texttt{\tiny{[GM$|$GM]}}
                \texttt{end game} \\
            }
        }
    }
    & & \\ \\

\end{supertabular}
}

\end{document}
